%%%%%%%%%%%%%%%%%%%%%%%%%%%%%%%%%%%%%%%%%%%%%%%%%%%%%%%%%%%%%%%%%%%%%%%%%%%%%%%%%%%%%%%%%%%%%%%%%%%%
%                                                                                                  %
%                     class options                                                                %
%                     options can be modified                                                      %
%                                                                                                  %
% -> usebiber - default option: uses the biblatex package                                          %
% -> usebibtex - use the classical bibtex with plain format                                        %
%%%%%%%%%%%%%%%%%%%%%%%%%%%%%%%%%%%%%%%%%%%%%%%%%%%%%%%%%%%%%%%%%%%%%%%%%%%%%%%%%%%%%%%%%%%%%%%%%%%%

\documentclass[usebibtex]{gammas}

%%%%%%%%%%%%%%%%%%%%%%%%%%%%%%%%%%%%%%%%%%%%%%%%%%%%%%%%%%%%%%%%%%%%%%%%%%%%%%%%%%%%%%%%%%%%%%%%%%%%
%                                                                                                  %
%                     author settings                                                              %
%                                                                                                  %
%%%%%%%%%%%%%%%%%%%%%%%%%%%%%%%%%%%%%%%%%%%%%%%%%%%%%%%%%%%%%%%%%%%%%%%%%%%%%%%%%%%%%%%%%%%%%%%%%%%%

\gammtitle{Article Template for GAMM Archive for Students}

% Please provide for each author the name and address (the full postal address must be provided in the online form but not on the paper)

% authors (there is a maximum of six authors possible)
% please mark the corresponding author with the command \coraauth after the institute (see the example below)
\gammauthora{Author one\inst{a}\corauth}
\gammauthorb{Author two\inst{a}}
\gammauthorc{Author three\inst{b}}

% authors are encouraged to include their orcid in the manuscript
% a circle placeholder will appear next to the author in the header
\gammauthoraorcid{1234-5678-1234-5678}

% Use the short forms of the author names for the header
\gammauthorhead{O.~Author, T.~Author, T.~Author}

% author affiliations (full postal address must be provided in the online form but not on the manuscript)
\gammaddressb{Technische Universität Berlin, Berlin, Germany} 
\gammaddressa{Technische Universität Chemnitz, Chemnitz, Germany}	% 

% additional information of the corresponding author
\gammcorrespondence{email@example.com}

% Supervisor: optional, if the supervisor wants to be mentioned
\gammsupervisor{FirstName LastName, TU Gammburg, Asymtopia}

% name of the bibfile without '.bib'
\gammbibfilename{example_bibliography}
\usepackage{blindtext}

%%%%%%%%%%%%%%%%%%%%%%%%%%%%%%%%%%%%%%%%%%%%%%%%%%%%%%%%%%%%%%%%%%%%%%%%%%%%%%%%%%%%%%%%%%%%%%%%%%%%
%	                  newcommands                                                                  %
%	                  this is optional but it makes live very easy                                 %
%%%%%%%%%%%%%%%%%%%%%%%%%%%%%%%%%%%%%%%%%%%%%%%%%%%%%%%%%%%%%%%%%%%%%%%%%%%%%%%%%%%%%%%%%%%%%%%%%%%%

% Theorem-like environments
\newtheorem{theorem}{Theorem}
\newtheorem{lemma}[theorem]{Lemma}

\theoremstyle{definition}
\newtheorem{definition}[theorem]{Definition}

%%%%%%%%%%%%%%%%%%%%%%%%%%%%%%%%%%%%%%%%%%%%%%%%%%%%%%%%%%%%%%%%%%%%%%%%%%%%%%%%%%%%%%%%%%%%%%%%%%%%
%                     begindocument start                                                          %
%                     do not modify this                                                           %
%%%%%%%%%%%%%%%%%%%%%%%%%%%%%%%%%%%%%%%%%%%%%%%%%%%%%%%%%%%%%%%%%%%%%%%%%%%%%%%%%%%%%%%%%%%%%%%%%%%%          
%
\begin{document}
\maketitle
%
%%%%%%%%%%%%%%%%%%%%%%%%%%%%%%%%%%%%%%%%%%%%%%%%%%%%%%%%%%%%%%%%%%%%%%%%%%%%%%%%%%%%%%%%%%%%%%%%%%%%
%                     begindocument end                                                            %
%                     do not modify this                                                           %
%%%%%%%%%%%%%%%%%%%%%%%%%%%%%%%%%%%%%%%%%%%%%%%%%%%%%%%%%%%%%%%%%%%%%%%%%%%%%%%%%%%%%%%%%%%%%%%%%%%%

%%%%%%%%%%%%%%%%%%%%%%%%%%%%%%%%%%%%%%%%%%%%%%%%%%%%%%%%%%%%%%%%%%%%%%%%%%%%%%%%%%%%%%%%%%%%%%%%%%%%
%                     abstract                                                                     %
%%%%%%%%%%%%%%%%%%%%%%%%%%%%%%%%%%%%%%%%%%%%%%%%%%%%%%%%%%%%%%%%%%%%%%%%%%%%%%%%%%%%%%%%%%%%%%%%%%%%
\begin{gammabstract}
	This is the abstract (max.\ 250 words for research papers and tutorials, max.\ 150 words for technical briefs)
\end{gammabstract}

%%%%%%%%%%%%%%%%%%%%%%%%%%%%%%%%%%%%%%%%%%%%%%%%%%%%%%%%%%%%%%%%%%%%%%%%%%%%%%%%%%%%%%%%%%%%%%%%%%%%
%                     keywords                                                                     %
%%%%%%%%%%%%%%%%%%%%%%%%%%%%%%%%%%%%%%%%%%%%%%%%%%%%%%%%%%%%%%%%%%%%%%%%%%%%%%%%%%%%%%%%%%%%%%%%%%%%
\begin{gammkeywords}
	mathematics, mechanics, \LaTeX \\ (max. 5 keywords)
\end{gammkeywords}

\section{Introduction}
This is the introduction. The authors are encouraged to include their \textsf{ORCiD} in the manuscript. The circle next to the author will be replaced by the logo during copyediting.

\section{Contents}
This is some scientific content. Make sure to use the commands from the \texttt{siunitx} package for all numbers. For instance, we have
\begin{gather*}
	\SI{1.2345e6}{\metre\per\second}, \qquad \SI{20}{\percent},\\
	{\sisetup{fixed-exponent=-9,scientific-notation=fixed}\SIrange{80e-9}{120e-9}{\kilogram}}, \qquad \num{5.4321e-9}.
\end{gather*}
For more details about the typesetting in GAMMAS, we refer to \cite{kern2019gammas}. For a concise introduction to scientific writing
in mathematics we refer to the excellent monograph~\cite{higham2020handbook}. Consult the documentation of the \texttt{amsmath}\footnote{\url{https://ctan.org/pkg/amsmath}} package for useful hints on typesetting equations in \LaTeX{}.

\section{Summary}
This is a great summary. 

%%%%%%%%%%%%%%%%%%%%%%%%%%%%%%%%%%%%%%%%%%%%%%%%%%%%%%%%%%%%%%%%%%%%%%%%%%%%%%%%%%%%%%%%%%%%%%%%%%%%
%                     code availability block (not optional)                                       %
%%%%%%%%%%%%%%%%%%%%%%%%%%%%%%%%%%%%%%%%%%%%%%%%%%%%%%%%%%%%%%%%%%%%%%%%%%%%%%%%%%%%%%%%%%%%%%%%%%%%
% Every accepted article for GAMMAS must contain a code availability statement. 
% You can modify the following example to account for your requirements.
\begin{gammacode}
	Both, the \LaTeX{} source code for reference~\cite{kern2019gammas} as well as the python-script for the simulation, are available as supplementary material and can be obtained under the
	\begin{center}
	 	\href{https://doi.org/10.14464/gammas.v1i1.417}{DOI:~10.14464/gammas.v1i1.417}.
	\end{center}
	% If you want to attach the code as supplementary material, then this DOI will be the same as the DOI for your article and inserted by the copyeditor. Otherwise, you can use another stable location.
	%
	If the source code is not shared, a simple statement should be given for explanation, e.g., for license issues, non-disclosure agreements or trade secrets. For \emph{best practices} for code we refer to \cite{fehr2016best}.
\end{gammacode}

%%%%%%%%%%%%%%%%%%%%%%%%%%%%%%%%%%%%%%%%%%%%%%%%%%%%%%%%%%%%%%%%%%%%%%%%%%%%%%%%%%%%%%%%%%%%%%%%%%%%
%                     acknowledgment block (optional)                                              %
%%%%%%%%%%%%%%%%%%%%%%%%%%%%%%%%%%%%%%%%%%%%%%%%%%%%%%%%%%%%%%%%%%%%%%%%%%%%%%%%%%%%%%%%%%%%%%%%%%%%
\begin{gammacknowledgement}
	If you received any funding for carrying out the reported research, you should list the funding agencies and (if applicable) the project and project number.
\end{gammacknowledgement}

%%%%%%%%%%%%%%%%%%%%%%%%%%%%%%%%%%%%%%%%%%%%%%%%%%%%%%%%%%%%%%%%%%%%%%%%%%%%%%%%%%%%%%%%%%%%%%%%%%%%
%                     appendix (optional)                                                          %
%%%%%%%%%%%%%%%%%%%%%%%%%%%%%%%%%%%%%%%%%%%%%%%%%%%%%%%%%%%%%%%%%%%%%%%%%%%%%%%%%%%%%%%%%%%%%%%%%%%%
\appendix

\section{Appendix}
This is an (optional) appendix.

%%%%%%%%%%%%%%%%%%%%%%%%%%%%%%%%%%%%%%%%%%%%%%%%%%%%%%%%%%%%%%%%%%%%%%%%%%%%%%%%%%%%%%%%%%%%%%%%%%%%
%                     enddocument start                                                            %
%                     do not modify this                                                           %
%%%%%%%%%%%%%%%%%%%%%%%%%%%%%%%%%%%%%%%%%%%%%%%%%%%%%%%%%%%%%%%%%%%%%%%%%%%%%%%%%%%%%%%%%%%%%%%%%%%%
\end{document}